\documentclass[10pt,a4paper]{amsart}
\usepackage[margin=19mm]{geometry}
\usepackage{amsmath,amssymb,mathtools}
\usepackage[scr=rsfs,cal=euler]{mathalfa}
\usepackage{xfrac}
\usepackage[unicode,hidelinks,bookmarks=false]{hyperref}
\newcommand{\ie}{i.e.\ }
\newcommand{\cf}{cf.\ }
\newcommand{\resp}{respectively\ }
\newcommand{\Subsection}{\subsection}
\providecommand{\nobreakdash}{\nobreak\hbox{-}}
\setlength{\emergencystretch}{2em}
\multlinegap0pt
\usepackage{xcolor}
\usepackage{amssymb}
\usepackage{bm}
\usepackage[shortlabels]{enumitem}
\newtheorem{lemma}{Lemma}
\newcommand\gen[1]{\left\langle#1\right\rangle}
\def\nsgp{\trianglelefteq}
\title{The growth of residually soluble groups}
\author{Sean Eberhard, Elena Maini}
\date{Source: arXiv:2511.07018v2 (CC BY 4.0)}
\begin{document}
\maketitle

\noindent\emph{Source and licence.} The growth of residually soluble groups. Version arXiv:2511.07018v2 (CC BY 4.0), distributed by arXiv under the Creative Commons Attribution 4.0 International licence, http://creativecommons.org/licenses/by/4.0/, as recorded in the arXiv licence metadata for this exact version and shown on the abstract page https://arxiv.org/abs/2511.07018v2. Full licence notice: \texttt{licenses/arxiv-2511.07018-CC-BY-4.0.txt}.\par\smallskip

\noindent\textit{Editorial scope.} Counted unit: the article's lemma lemma (source0.tex lines 309--313). The article's complete original proof of it is delivered with it (source0.tex lines 314--351, one \texttt{proof} environment of 38 lines). No mathematics is written, repaired or re-derived: the statement and the proof are the article's own lines, in the article's own notation, and no retained line is substituted or re-flowed.  No further source-local context is needed: every cross-reference of every retained line resolves inside the two delivered spans.  Nothing was omitted from any retained span, so the package carries no omission record.  No retained line contains a citation, so the wrapper carries no bibliography and no bibliography entry.  No retained line contains a figure environment, an image-inclusion command or drawing code, so no image is delivered and the package declares no asset dependency.  Editorial formatting only, in addition: an AMS \texttt{amsart} preamble that restates the article's own declarations and macros used by the retained lines, namely the environments \texttt{lemma} on separate counters, together with the standard packages those lines need.  The retained environments print in this document as Lemma 1 in order of appearance, whereas the article prints the same items with the numbering of its own full document.  The complete native source of the article as distributed in the arXiv e-print for this exact version is bundled unchanged as \texttt{sources/188778/source0.tex}.
\begin{lemma}[Generalized Milnor lemma]
    \label[lemma]{lem:soft-generalized-milnor}
    Let $G = \gen X$ be a finitely generated group and assume that $H \le G$ is a finitely generated subgroup whose normal closure $N = \gen{H^G}$ is not finitely generated.
    Then $G$ has growth \[\gamma_X(n) \succeq \exp(n^{1/2}).\]
\end{lemma}

\begin{proof}
    Suppose $H = \gen{Y}$ where $Y$ is finite.
    Let
    \begin{equation}
        \label{eq:Y_i-def}
        Y_i = \{y^g : y \in Y,~\ell_X(g) \le i\}, \qquad H_i = \gen{Y_i} \qquad (i \ge 0).
    \end{equation}
    Then $H_0 \le H_1 \le \cdots$ is an increasing chain of subgroups of $N$. If $H_k = H_{k+1}$ for some $k$ then for $x \in X \cup X^{-1}$ we have
    \[
        H_k^x = \gen{Y_k^x} \le \gen{Y_{k+1}} = H_{k+1} = H_k,
    \]
    so $H_k \nsgp G$, so $H_k = N$.
    Since $N$ is not finitely generated, there is no such $k$.

    Therefore, for each $i > 0$ there is some $y_i \in Y_i \setminus H_{i-1}$.
    Now for all $k \ge 1$ we claim that the $2^k$ elements $y_1^{\epsilon_1} \cdots y_k^{\epsilon_k}$ with $\epsilon_i \in \{0,1\}$ are distinct. Indeed, suppose
    \[
        y_1^{\epsilon_1} \cdots y_k^{\epsilon_k} = y_1^{\delta_1} \cdots y_k^{\delta_k}
    \]
    where $\epsilon,\delta \in \{0,1\}^k$. Then
    \(
        H_{k-1} y_k^{\epsilon_k} = H_{k-1} y_k^{\delta_k},
    \)
    which implies $\epsilon_k = \delta_k$ since $y_k \notin H_{k-1}$.
    Cancelling $y_k^{\epsilon_k}$ and applying induction, it follows that $\epsilon = \delta$.

    Now observe that
    \[
        \ell_X(y_1^{\epsilon_1} \cdots y_k^{\epsilon_k}) \le \sum_{i=1}^k \ell_X(y_i) \le k (L + 2k),
    \]
    where $L = \ell_X(Y)$.
    Thus
    \(
        \gamma_X(kL + 2k^2) \ge 2^k
    \)
    for all $k \ge 1$,
    which implies that $\gamma_X(n) \succeq \exp(n^{1/2})$.
\end{proof}

\end{document}
